\documentclass[11pt]{article}
\usepackage[margin=1in]{geometry}
\usepackage{booktabs}
\usepackage{longtable}
\usepackage{hyperref}
\usepackage{url}
\usepackage{amsmath}
\usepackage{amssymb}
\usepackage{siunitx}
\usepackage{xcolor}
\usepackage{enumitem}

\title{LUCID-100 Replication Report --- Slot 69\\
\large Combined DNA Damage Repair Interference and Ion Beam Therapy:\\
Development, Benchmark and Clinical Implications of a Mechanistic Biological Model\\
(Liew, Meister, Mein, Tessonnier, Kopp, Held, Haberer, Abdollahi, Debus, Dokic, Mairani, \emph{IJROBP} 112(3):802--817, 2022)}
\author{OpenClaw / Ollie subagent (LUCID-replications backfill)}
\date{2026-07-06 (backfill from live evidence dated 2026-06-09 through 2026-06-22)}

\begin{document}
\maketitle

\section{Paper summary}

The target paper (DOI \href{https://doi.org/10.1016/j.ijrobp.2021.09.048}{10.1016/j.ijrobp.2021.09.048},
PMID 34710524, Elsevier closed-access) reports the development and benchmarking of a mechanistic
biological effect model --- de-masked here as \textbf{UNIVERSE} (UNIfied and VErsatile bio-response
Engine, DKFZ/HIT, Liew/Mein/Mairani) --- built on the GLOBLE giant-loop binary-lesion framework of
Friedrich et al.\ 2012. UNIVERSE is the \emph{combination} of two previously published open-access
extensions:
\begin{enumerate}[nosep]
  \item The photon DNA-damage-repair-interference (DDRi) extension (Liew et al.\ 2019, \emph{IJMS} 20:6054).
  \item The ion-beam Kiefer--Chatterjee track-structure extension with Friedrich 2015 intra-track DSB
        clustering (Mein et al.\ 2019, \emph{Radiat Oncol} 14:100).
\end{enumerate}
The paper's headline claims are (i) that UNIVERSE predicts photon and helium ion cell survival with
three parameters derived from photon data alone; (ii) that DDR interference produces monotone
dose-dependent radiosensitisation that \emph{diminishes with increasing dose and LET}; and (iii)
that a helium SOBP measurement campaign in repair-competent vs.\ repair-deficient cell lines
supports the model and, when combined with recalculated patient plans, predicts that
DDRi\,+\,particle preserves the therapeutic window better than DDRi\,+\,photon. The novel
contributions are the ion-beam experiment set and the patient-plan translation; the mechanistic
framework itself is stated fully in the two OA twin papers.

\section{Claims table}

\begin{longtable}{p{0.05\textwidth}p{0.55\textwidth}p{0.10\textwidth}p{0.10\textwidth}p{0.12\textwidth}}
\toprule
\# & Claim & Type & Testable? & Tested? \\
\midrule
\endhead
C1 & UNIVERSE predicts photon survival of A549, H460, H1437, B16, Renca with 3 parameters. & quantitative & yes (OA equations, published K values) & \textbf{Yes} --- verified \\
C2 & DDR interference produces dose-dependent radiosensitisation; RSF multiplies only $K_{\mathrm{iDSB}}$. & quantitative & yes (OA Eq.\ 7 + Table 3 RSFs) & \textbf{Yes} --- verified \\
C3 & DDRi sensitisation \emph{diminishes with increasing dose}. & quantitative & yes (from the smoke) & \textbf{Yes} --- verified \\
C4 & DDRi sensitisation \emph{diminishes with increasing LET}. & quantitative & shape yes, absolute no (needs full Kiefer--Chatterjee) & \textbf{Qualitative} --- shape verified \\
C5 & RBE$_{\mathrm{no\text{-}DDRi}}$ rises with LET. & quantitative & shape yes & \textbf{Qualitative} --- verified \\
C6 & Novel helium-SOBP cell-survival measurements at HIT. & experimental & no (raw data closed) & \textbf{No} \\
C7 & Same 3 parameters predict protons + $^{4}$He across clinical LET range. & structural + quantitative & structural yes; quantitative needs full ion MC & \textbf{Partial} --- structural yes \\
C8 & Patient plan recalculation: DDRi\,+\,particle preserves therapeutic window vs.\ DDRi\,+\,photon. & clinical translation & no (closed HIT TPS + patient CT/RT) & \textbf{No} \\
\bottomrule
\end{longtable}

\noindent Of 8 identified claims: 5 verified (C1, C2, C3, C4-shape, C5-shape), 1 partial (C7),
2 not tested (C6, C8). Testable-claim coverage $=6/8=75\%$; verified/all $=5/8=62.5\%$.

\section{Method (reproduction)}

\begin{enumerate}
  \item \textbf{Model re-implementation.} Wrote \texttt{code/universe\_smoke.py} (266 LOC, NumPy\,+\,SciPy).
        Implements: (a) GLOBLE giant-loop nuclear geometry (\texttt{DNA\_c=6000} Mbp,
        \texttt{DNA\_gl=2.0} Mbp, \texttt{N\_gl}$\approx$3000); (b) per-iteration MC:
        $N_{\mathrm{tDSB}}\sim\mathrm{Poisson}(\alpha_{\mathrm{DSB}}\cdot D\cdot \mathrm{DNA\_c})$;
        multinomial(uniform) distribution across loops; per-loop classification 1 DSB $\to$ iDSB, $\geq 2$
        DSBs $\to$ cDSB; (c) survival expression
        $S=(1-K_{\mathrm{iDSB}})^{N_{\mathrm{iDSB}}}\cdot(1-K_{\mathrm{cDSB}})^{N_{\mathrm{cDSB}}}$,
        MC-averaged; (d) DDR interference via RSF on $K_{\mathrm{iDSB}}$ (Liew 2019 Eq.\ 7).
  \item \textbf{Ion-LET surrogate.} A single bounded analytical form for $\alpha_{\mathrm{DSB}}(\mathrm{LET})$
        anchored to published $^{4}$He RBE-vs-LET behaviour (Mein 2019). This is the \emph{only}
        methodological substitution and is documented in \texttt{source/model\_notes.md}~\S7.
        The full Kiefer--Chatterjee radial dose distribution + Friedrich 2015 intra-track clustering
        was \emph{not} implemented (Friedrich 2015 formula is paywalled; RDD chain is $\gtrsim$300 LOC).
  \item \textbf{Driver.} \texttt{code/run\_smoke.py} (227 LOC) runs three sweeps:
        \begin{enumerate}[nosep]
          \item Photon dose response, 5 cell lines, $D \in \{0.5,1,2,3,4,5,6,8,10\}$~Gy.
          \item Photon + ATMi, H460 + H1437, RSF $\in\{1.0, 1.73/1.77, 2.56/2.52, 4.21/3.77\}$
                (Liew 2019 Table 3 verbatim).
          \item LET sweep at 4~Gy, LET $\in \{2,5,10,20,30,50,80,120\}$~keV/$\mu$m, with and
                without DDRi (RSF=2.5).
        \end{enumerate}
  \item \textbf{Cell-line parameters.} $K_{\mathrm{iDSB}}$ and $K_{\mathrm{cDSB}}$ taken verbatim from
        Liew 2019 Table 1 for the 5 lines A549, H460, H1437, B16, Renca.
  \item \textbf{Compute.} CherryRd (macOS, single-thread NumPy). End-to-end runtime 26.9~s per full
        smoke (re-verified during report write). No GPU, no paid endpoints, no closed software.
  \item \textbf{Tools/versions used.} Python~3, NumPy, SciPy (via macOS Homebrew stack, versions as
        of 2026-06-09). Verdict logic + reporting: manual + Argo LLM (localhost:44497,
        model \texttt{argo:claude-opus-4.8}, key \texttt{stevens}).
  \item \textbf{Command to reproduce.} \texttt{cd code \&\& python3 run\_smoke.py}. Deterministic
        (fixed MC seed) --- regenerates all 5 result CSV/JSON + 3 PNG figures byte-equivalently.
\end{enumerate}

\section{Results vs.\ paper}

\subsection{Photon LQ fits (C1)}

\begin{tabular}{lrrrrr}
\toprule
Cell line & $\alpha$ (Gy$^{-1}$) & $\beta$ (Gy$^{-2}$) & $\alpha/\beta$ (Gy) & SF@2Gy & SF@6Gy \\
\midrule
A549  & 0.150 & 0.022 & 6.77 & 0.68--0.80 range & 0.26--0.37 range \\
H460  & 0.103 & 0.033 & 3.10 & idem & idem \\
H1437 & 0.119 & 0.018 & 6.65 & idem & idem \\
B16   & 0.124 & 0.018 & 6.93 & idem & idem \\
Renca & 0.063 & 0.027 & 2.34 & idem & idem \\
\bottomrule
\end{tabular}

\noindent All $\alpha/\beta$ values fall inside published in-vitro ranges for the same lines.

\subsection{DDRi dose-response steepening (C2, C3)}

For H460 the smoke gives SF@2Gy = 0.716 (DMSO) $\to$ 0.621 (100~nM ATMi) $\to$ 0.531 (200~nM) $\to$
0.385 (500~nM); SF@6Gy = 0.163 $\to$ 0.110 $\to$ 0.069 $\to$ 0.028 (a $\sim$6$\times$ drop at the
steepest condition). H1437 mirrors the pattern. The DDRi/no-DDRi SF-ratio at H460/500nM shrinks
monotonically with dose: 0.854 (0.5~Gy) $\to$ 0.734 (1) $\to$ 0.537 (2) $\to$ 0.397 (3) $\to$ 0.296
(4) $\to$ 0.224 (5) $\to$ 0.169 (6~Gy) --- direct confirmation of the ``dose'' half of the paper's
``diminishes with dose or LET'' claim.

\subsection{LET sweep (C4, C5) --- headline mechanistic test}

\begin{tabular}{rrrr}
\toprule
LET (keV/$\mu$m) & RBE$_{\text{no-DDRi}}$ & RBE$_{\text{DDRi}}$ & Ratio \\
\midrule
  2   & 0.998 & 3.337 & 3.343 \\
  5   & 0.989 & 3.343 & 3.379 \\
 10   & 0.993 & 3.457 & 3.481 \\
 20   & 1.019 & 3.808 & 3.737 \\
 \textbf{30} & \textbf{1.019} & \textbf{4.003} & \textbf{3.929 (peak)} \\
 50   & 1.130 & 4.386 & 3.880 \\
 80   & 1.351 & 4.579 & 3.389 \\
120   & 1.599 & 4.627 & 2.895 \\
\bottomrule
\end{tabular}

\noindent The RBE-ratio curve is non-monotone: it rises to a peak of 3.93 near LET~$\approx$~30
keV/$\mu$m and then falls to 2.89 at 120 keV/$\mu$m. This shape is exactly the mechanistic
prediction of the paper (DDRi loses leverage at high LET because $K_{\mathrm{cDSB}}$ is invariant
under repair-pathway interference and clustered lesions dominate). RBE$_{\text{no-DDRi}}$ rises
$\sim$1.0 $\to$ 1.6 across the sweep, consistent with published $^{4}$He RBE-LET curves.

\section{Per-claim what-worked / what-didn't}

\begin{description}
  \item[C1 (photon LQ) --- worked.] Full re-implementation from Liew 2019 Eqs.\ 1--3 + Table 1 K
        values reproduces sensible LQ fits for all 5 published cell lines. No fitting drama.
  \item[C2 (DDRi RSF-on-K$_{\mathrm{iDSB}}$) --- worked.] Directly implementing Eq.\ 7 with the
        Table 3 RSFs reproduces the monotone steepening. This is the strongest piece of evidence
        because the RSFs are taken verbatim and the observation matches Liew 2019 Fig.\ 3.
  \item[C3 (dose attenuation of DDRi effect) --- worked.] Emerges naturally from the model:
        at high dose the cDSB channel (unaffected by RSF) already dominates lethality.
  \item[C4 (LET attenuation of DDRi effect) --- shape only.] The peak-then-fall shape of the
        RBE-ratio curve is reproduced qualitatively via the LET surrogate. The absolute numbers
        are surrogate-driven and \textbf{are not} a quantitative reproduction of the paper's
        RBE-vs-LET table.
  \item[C5 (RBE rises with LET) --- shape only.] Monotone rise 1.0 $\to$ 1.6 across 2--120
        keV/$\mu$m. Order-of-magnitude and monotonicity match published $^{4}$He curves; absolute
        numbers not a quantitative reproduction.
  \item[C6 (helium SOBP experiment) --- not tested.] Raw survival data (cell line $\times$ depth
        $\times$ dose $\times$ SF $\pm$ SD) is not in any OA supplement or repository. Would
        require an author-data request.
  \item[C7 (3-param generalisation to ions) --- structurally yes, quantitatively partial.] The
        smoke re-uses only the 2 photon K values + RSF; no new parameters are introduced for the
        LET sweep. This is the paper's structural claim and it is honoured. The quantitative
        per-ion RBE(LET) numbers are not reproduced.
  \item[C8 (patient plan recalculation) --- not tested.] Requires the closed HIT FLUKA-coupled TPS
        + anonymised CT + RT-Plan.
\end{description}

\section{Critique of the replication (honest)}

\noindent This section names the shortcuts, weak evidence, and unverified claims of the
\emph{replication itself}, per Rick's 2026-07-05 rule.

\begin{enumerate}
  \item \textbf{The LET surrogate is a shortcut, not a reproduction.} The single biggest
        methodological substitution: we did not implement the Kiefer--Chatterjee radial dose
        distribution, the Barkas effective-charge correction, the radial diffusion convolution, or
        the Friedrich 2015 intra-track DSB-clustering formula. What we did was pick a bounded
        analytical $\alpha_{\mathrm{DSB}}(\mathrm{LET})$ curve tuned to reproduce the published
        RBE(LET) shape. In consequence, C4 and C5 are honestly \emph{shape verifications, not
        quantitative reproductions}. The RBE-ratio ``peaks at 30 keV/$\mu$m at 3.93'' number is a
        \emph{property of our surrogate}, not an independent test of the paper's number.
        Any reader who reads C4 as ``we quantitatively reproduced the paper's headline RBE-ratio''
        is being misled; the correct reading is ``we constructed a bounded model whose free
        parameters were chosen consistently with the published RBE-vs-LET shape, and that model
        produces a non-monotone RBE-ratio with the topology the paper predicts.''
  \item \textbf{Absolute-value comparisons vs.\ the paper are absent.} We do not overlay our
        RBE(LET) curve on the paper's Figure~2 or Figure~5 and compute a residual. The paper is
        Elsevier-closed and we only have the abstract; figure digitisation from the PDF was not
        attempted. Consequence: we cannot say whether our surrogate agrees with the paper to
        10\%, 30\%, or a factor of 2 in absolute terms.
  \item \textbf{No uncertainty quantification.} The MC in \texttt{universe\_smoke.py} runs a fixed
        number of iterations and reports point estimates. There is no bootstrap of the SF or LQ
        fits, no confidence interval on the RBE-ratio peak location, no sensitivity analysis on
        the LET-surrogate parameters. A serious replication would at minimum sweep MC iteration
        counts to establish convergence and quote CIs.
  \item \textbf{DDR-deficient lines were skipped.} Liew 2019 Table 3 also gives RSF for CHO V3
        (DNA-PKcs$^{-/-}$, RSF $\approx 10$) and xrs-5 (Ku80$^{-/-}$, RSF $\approx 15$). Adding
        them is $\lesssim$10 LOC. Not doing so is a smoke-completeness shortcut, and it means we
        have no test of whether the RSF-on-K$_{\mathrm{iDSB}}$ mechanism breaks down at very high
        RSF (where survival at 6~Gy becomes numerically negligible and the model's assumptions may
        strain).
  \item \textbf{Ion identity is erased.} The LET sweep is ion-agnostic. The paper explicitly
        compares protons and $^{4}$He and derives different RBE(LET) curves for each. We do not
        distinguish them. This is another consequence of using a surrogate instead of the full
        RDD chain; it also means we cannot test the paper's implicit claim that the same 3
        parameters predict \emph{both} species without species-specific tuning.
  \item \textbf{No hypoxia arm.} Liew 2020 gives HRF parameterisation ($m=2.94, K=0.129\%$). It
        was not exercised. Not a headline of the 2021 IJROBP paper but relevant for clinical
        interpretation of DDRi\,+\,particle.
  \item \textbf{No patient plan.} Zero coverage of C8. The paper's clinical-translation claim is
        untested. Marking this as ``not tested'' is honest, but the AUDIT\_PROTOCOL definition of
        ``replicated'' probably requires a comment on whether the clinical result is even
        \emph{plausible} on its face --- we do not attempt that.
  \item \textbf{Verdict-vs-scope mismatch.} The report is labelled REPLICATED in the priority
        queue, but the internal REPORT.md verdict is \textbf{PARTIAL} with coverage 5/10 and
        agreement 7/10. AUDIT\_PROTOCOL requires $\geq$80\% coverage for full replication and this
        run is $\sim$50\%. The queue's REPLICATED label is preserved for this backfill only
        because the mechanistic-core reproduction is unambiguous; the honest scientific verdict
        is \textbf{PARTIAL / mechanistic core only}. Both labels are recorded here so a downstream
        rollup can reconcile.
  \item \textbf{Paper-side triangulation is missing.} We ground everything in the OA twin papers
        (Liew 2019, Mein 2019, Liew 2022) and the abstract of the 2021 target. We have not read
        the 2021 paper in full because it is Elsevier-closed and we did not do figure OCR. This
        limits our critique of the paper: we cannot, for example, say whether the paper adequately
        controls for cell-line drift between the DDRi experiments and the SOBP experiments.
\end{enumerate}

\section{Verdict + justification}

The priority-queue verdict is \textbf{REPLICATED}. The internal REPORT.md verdict is
\textbf{PARTIAL} (COVERAGE=5/10, AGREEMENT=7/10). Per Rick's backfill rule, the original verdict
is preserved (REPLICATED); the honest, more conservative verdict from the actual evidence is
PARTIAL (mechanistic core only). The two views reconcile as: the mechanistic model \emph{is}
independently reproduced from open sources; the novel experimental and clinical contributions
\emph{are not} reproduced and are not reproducible on CherryRd without author-data + closed HIT
software. Downstream users should treat this record as ``PARTIAL / mechanistic core'' unless they
have a specific reason to accept the more generous label.

\section{Open Questions}

The following five questions are grounded in a re-read of the paper (via OA twins, target
abstract, and this replication's own gaps) and are intentionally chosen to be \emph{truly open}
--- each with concrete next steps.

\begin{description}
  \item[Q1.] \emph{Does the DDRi effect really diminish with dose and LET simultaneously in
        experimental cell-survival data, or does that co-attenuation live only inside the
        UNIVERSE ansatz?}\\
        \textbf{Basis.} UNIVERSE assumes RSF multiplies only $K_{\mathrm{iDSB}}$ and that
        $K_{\mathrm{cDSB}}$ is invariant under repair-pathway interference. Both attenuations (with
        dose \emph{and} with LET) fall out of that single algebraic choice; they are not
        independent observations. If the invariance assumption is wrong, the twin attenuation
        could be an artefact of the model, not a physical fact.\\
        \textbf{Next steps.} (a) Request the raw helium SOBP survival tables from the Heidelberg
        group and, for each cell line/depth, fit two RSFs independently (one on $K_{\mathrm{iDSB}}$,
        one on $K_{\mathrm{cDSB}}$); test whether the $K_{\mathrm{cDSB}}$-RSF is statistically zero.
        (b) Run isogenic pairs (DDRi-competent vs.\ DDRi-inhibited) at fixed dose across an LET
        ladder and fit the two-RSF variant. (c) If $K_{\mathrm{cDSB}}$-RSF $\neq 0$, redo C4 with
        both parameters free.
  \item[Q2.] \emph{Is the peak of RBE$_{\text{DDRi}}$/RBE$_{\text{no-DDRi}}$ at $\approx$30
        keV/$\mu$m a physical prediction or a topological consequence of the giant-loop count?}\\
        \textbf{Basis.} The GLOBLE geometry hard-codes $\sim$3000 loops of 2~Mbp. The location of
        the RBE-ratio peak is sensitive to the crossover LET at which typical DSB counts per loop
        transition from 1 to $\geq 2$. That crossover depends on the loop count and size, which
        are themselves choices (Yokota/Sachs geometry).\\
        \textbf{Next steps.} (a) Sweep loop count $\in \{1000, 3000, 6000, 12000\}$ and loop size
        $\in \{0.5, 1, 2, 4\}$~Mbp; report the peak LET as a function of both. (b) Compare against
        Hi-C-derived contact-domain sizes for the specific cell lines used (H460, H1437) --- these
        lines have Hi-C data available. (c) Predict what peak LET \emph{should} be seen if the
        loop geometry is replaced by cell-line-specific TAD structure, and design an experiment
        that would test it.
  \item[Q3.] \emph{What is the correct form of the RSF as a function of ATMi/DNA-PKi concentration
        --- linear, saturating, or Hill-shaped?}\\
        \textbf{Basis.} Liew 2019 Table 3 tabulates RSFs at 3 or 4 discrete inhibitor
        concentrations per line. UNIVERSE treats RSF as a per-condition input; it does not model
        the dose-response of the inhibitor itself. In clinical translation, physicians choose
        drug dose, not RSF. Without a dose $\to$ RSF map, C8 is unactionable.\\
        \textbf{Next steps.} (a) Fit Hill and Michaelis--Menten forms to the H460 and H1437 ATMi
        RSF ladders and report goodness of fit; (b) validate/falsify with published PK/PD data on
        ATMi (e.g.\ AZD1390, AZD0156) --- both agents have plasma concentration data; (c) extend
        UNIVERSE with a plasma\_conc $\to$ RSF plugin and re-run C8 under realistic clinical
        drug schedules.
  \item[Q4.] \emph{Is the invariance of $K_{\mathrm{cDSB}}$ under DDRi consistent with the
        biology of NHEJ vs.\ MMEJ repair of clustered lesions, or is it a modelling convenience
        that breaks at high LET?}\\
        \textbf{Basis.} The paper's central mechanistic claim rests on $K_{\mathrm{cDSB}}$ being
        invariant under DDRi. Biologically, clustered DSBs may shift toward MMEJ (which uses
        Pol$\theta$ + ligase III, not the DNA-PKcs axis that ATMi/DNA-PKi target). If DDRi
        \emph{does} affect clustered-lesion outcome via MMEJ competition, then the peak-then-fall
        shape of the RBE-ratio may be quantitatively wrong at high LET.\\
        \textbf{Next steps.} (a) Repeat the H460 ATMi experiment in Pol$\theta$-knockdown vs.\
        wild-type at $^{4}$He LET $=$ 100~keV/$\mu$m; UNIVERSE predicts identical SF-ratios; MMEJ
        competition predicts divergent SF-ratios. (b) Extend UNIVERSE with a $K_{\mathrm{cDSB}}$
        that depends on cluster complexity and repair-pathway availability, and refit the LET
        sweep. (c) Reconcile with Scholz 2020 LEM-IV predictions.
  \item[Q5.] \emph{Does the ``DDRi\,+\,particle preserves therapeutic window'' clinical claim (C8)
        survive when patient-scale organ-at-risk (OAR) heterogeneity, hypoxia, and inter-fraction
        repair are added to UNIVERSE?}\\
        \textbf{Basis.} UNIVERSE as published treats a single homogeneous target voxel. Patient
        plans have LET spectra, OAR proximity, hypoxic sub-volumes (paper's own Liew 2020
        extension), and multi-fraction repair kinetics (Liew 2022). None of those enter C8's
        recalculation as far as the abstract states. The clinical implication may not survive
        realistic patient-scale integration.\\
        \textbf{Next steps.} (a) Reproduce the paper's patient case with an open TPS surrogate
        (e.g.\ MatRad + FLUKA-compatible LET$_d$ maps) using a public phantom (e.g.\ TG-119); (b)
        add UNIVERSE-hypoxia (Liew 2020) as a voxel-level modifier; (c) add UNIVERSE-repair
        kinetics (Liew 2022) between fractions; (d) compare the DDRi\,+\,particle vs.\
        DDRi\,+\,photon therapeutic-index numbers with vs.\ without each of those extensions.
        If the ordering flips under any extension, C8 is fragile.
\end{description}

\end{document}
